\section{New restrictions on proposed split-eightfold CM
inputs}\label{new-restrictions-on-proposed-split-eightfold-cm-inputs}

Author/checker: GPT-6 Astra Pro; inherited effort not exposed;
HODGE\_DIVISION /root/sixfolds. Date: 2026-11-21.

These results address the remaining nonsplit sixfold gate through a
possible split-eightfold supplier. They take DEG\_L S-21--S-25 and
S-30--S-34 as established background and do not repeat the excluded
genus-four ample imaginary-secant square. They do not construct the
missing eightfold. The first restriction excludes direct-sum repairs of
the imaginary-secant character. The second excludes a different natural
source: the fully semiregular product of two surface point ideals, even
after derived equivalence, isogeny and rational divisor normalization.

\subsection{SF-E801. Intrinsic exterior
kernels}\label{SF-E801}

Let \(X\) be an abelian fourfold and put

\[
 W_X=H^0(X,T_X)\oplus H^1(X,\mathcal O_X)
     =\operatorname{Lie}(X\times\widehat X).
\]

The Hochschild cohomology algebra is the exterior algebra of \(W_X\).
Choose a nonzero holomorphic four-form \(\Omega\). Contraction by the
first summand and exterior multiplication by the second identify the
full Hodge cohomology with the free exterior module generated by
\(\Omega\). In particular, a diagonal Hodge character
\(v\in\bigoplus_pH^p(\Omega_X^p)\) corresponds to a four-vector
\(\widetilde v\in\bigwedge^4W_X\), and its degree-two contraction map
becomes

\[
 c_v^2:\bigwedge^2W_X\longrightarrow\bigwedge^6W_X,
 \qquad z\longmapsto z\wedge\widetilde v.
 \tag{E801.1}
\]

Changing \(\Omega\) multiplies every displayed four-vector by a nonzero
scalar and changes none of the kernels below.

\subsubsection{The secant kernel}\label{the-secant-kernel}

Let \(W=L\oplus M\), with both summands of dimension four, and let

\[
 v=\lambda\operatorname{vol}_L+\mu\operatorname{vol}_M,
 \qquad\lambda\mu\ne0.
\]

Then

\[
 \ker(z\mapsto z\wedge v)=L\otimes M,
 \qquad\operatorname{rank}c_v^2=12.
 \tag{E801.2}
\]

Indeed, the two pure exterior-square summands map injectively into
distinct bidegrees, while the mixed summand annihilates both volume
forms. Moreover,

\[
 \{w\in\bigwedge^4W:(L\otimes M)\wedge w=0\}
       =\bigwedge^4L\oplus\bigwedge^4M.
 \tag{E801.3}
\]

To prove this last assertion, decompose \(w\) by its \(L,M\) bidegree
((p,4-p)). For \(1\le p\le3\), exterior multiplication gives an
injective map

\[
 \bigwedge^pL\longrightarrow
       \operatorname{Hom}(L,\bigwedge^{p+1}L).
\]

The corresponding map for \(M\) is also injective. Their tensor product
says that a nonzero component of bidegree ((p,4-p)) cannot be
annihilated by every mixed bivector. The two extreme components are
annihilated. This proves (E801.3).

\subsubsection{The product kernel and recovery of its two
factors}\label{the-product-kernel-and-recovery-of-its-two-factors}

Suppose instead that \(W=U_1\oplus U_2\), each \(U_i\) has dimension
four, and \(\omega_i\in\bigwedge^2U_i\) is nondegenerate. Put
\(v=\omega_1\wedge\omega_2\). Then

\[
 \ker c_v^2=K_1\oplus K_2,
 \qquad K_i=\{z\in\bigwedge^2U_i:z\wedge\omega_i=0\},
 \qquad\dim K_i=5,
 \tag{E801.4}
\]

and \(c_v^2\) has rank eighteen. On the two pure exterior-square
summands its image has dimension one each, in bidegrees ((4,2)) and
((2,4)). On \(U_1\otimes U_2\), it is the tensor product of the
isomorphisms \(u\mapsto u\wedge\omega_i\), so has rank sixteen in
bidegree ((3,3)).

The decomposable nonzero bivectors in the projective kernel form exactly
two disjoint smooth quadric threefolds, one in each \(\mathbf P(K_i)\).
To see that no mixed point occurs, write \(z=z_1+z_2\in K_1\oplus K_2\).
If \(z\) is decomposable, then \(z^2=0\). Its bidegree ((2,2)) part is
\(2z_1\otimes z_2\), forcing one \(z_i\) to vanish. Within
\(\bigwedge^2U_i\), decomposability is the Plücker quadratic equation
\(z_i^2=0\). Its restriction to the hyperplane \(K_i=\omega_i^\perp\) is
nondegenerate, because the exterior-square pairing is nondegenerate and
\(\omega_i^2\ne0\). It is therefore a smooth irreducible quadric in
\(\mathbf P^4\).

The two original four-dimensional summands are recovered intrinsically
from these two components: for each component, take the span in \(W\) of
the two-planes represented by its points. This span is \(U_i\). Every
nonzero vector in a symplectic four-space lies in a Lagrangian
two-plane, which verifies the last claim directly. Thus the unordered
pair \({U_1,U_2}\) is determined by \(ker c_v^2\), without a chosen
presentation of \(v\).

The accompanying exact script verifies the ranks (12,18), the
two-dimensional common annihilator in (E801.3), and the rank-five
Plücker quadratic forms. The geometric component and recovery arguments
are the proofs above, not inferences from numerical sampling.

One further common-annihilator identity is useful:

\[
 \{w\in\bigwedge^4W:(K_1\oplus K_2)\wedge w=0\}
       =\mathbf C(\omega_1\wedge\omega_2).
 \tag{E801.5}
\]

For a bidegree \((p,4-p)\) component, the equations for \(K_1\) kill
\(p=0,1\), and those for \(K_2\) kill \(p=3,4\). At \(p=1\), the
bivectors annihilating a fixed nonzero vector form a three-dimensional
space and cannot contain the five-dimensional \(K_1\); applying
coefficient functionals gives the tensor-valued assertion. At \(p=2\),
orthogonality to the primitive hyperplane makes the first factor
proportional to \(\omega_1\), and the second condition makes the second
factor proportional to \(\omega_2\). This proves (E801.5), also checked
by the exact script.

\subsection{SF-E802. Direct sums cannot repair the imaginary-secant
character}\label{SF-E802}

Let \(E\) be a gluable fully semiregular perfect object on an abelian
fourfold. Suppose that, for a rational algebraic divisor class \(B\), a
rational ample class \(\theta\), and \(s^2=-d<0\),

\[
 e^B\operatorname{ch}(E)
 =a\frac{e^{s\theta}+e^{-s\theta}}2
  +b\frac{e^{s\theta}-e^{-s\theta}}{2s},
 \qquad (a,b)\in\mathbf Q^2\setminus\{(0,0)\}.
 \tag{E802.1}
\]

In any decomposition \(E=\bigoplus_jE_j\), at most one summand has
nonzero rational Chern character. In particular, if \(E\) is a coherent
sheaf, it is indecomposable. The same assertion holds for a
decomposition in an effective finite equivariant category of the kind in
DEG\_L S-21.

\textbf{Proof.} The B-field acts by an invertible shear of \(W_X\), so
it can be applied to all the characters at once. The two exponential
characters become decomposable four-vectors with transverse
four-dimensional supports. Transversality follows from nondegeneracy of
\(2s\theta\). Their coefficients \((a\pm b/s)/2\) are both nonzero,
since \(a,b\) are rational and \(s\) is nonreal. Thus (E801.2) applies.

The CY4 adjunction used in S-21 says that injectivity of the full trace
\(\sigma_E^2\) makes Hochschild evaluation onto, so

\[
 \operatorname{im}\sigma_E^2=\operatorname{im}c_E^2,
 \qquad\dim\operatorname{Ext}^2(E,E)=12.
\]

For a direct sum, the Atiyah class is block diagonal. The trace is zero
on off-diagonal self-Ext blocks and is the usual full trace on each
diagonal block. Injectivity therefore makes the diagonal trace images
linearly independent. For each summand,

\[
 \operatorname{im}c_{E_j}^2\subset
 \operatorname{im}\sigma_{E_j}^2\subset\operatorname{im}c_E^2.
\]

The wedge map in (E801.1), paired with another bivector to obtain a top
form, is a symmetric bilinear form on \(\bigwedge^2W_X\). Hence the
image inclusion implies

\[
 \ker c_E^2\subset\ker c_{E_j}^2.
\]

Equation (E801.3) now forces each normalized summand character into the
same two-dimensional secant plane. Rationality makes it
\(a_j\alpha+b_j\beta\) with \(a_j,b_j\in\mathbf Q\): the rank determines
\(a_j\), and the degree-two cohomology component determines
\(b_j\theta\). If this character is nonzero, both pure coefficients are
again nonzero, so its degree-two contraction rank is twelve. Two such
summands would give two independent trace images of dimension at least
twelve inside a twelve-dimensional space. This is impossible.

Every nonzero coherent sheaf has nonzero rational Chern character: its
first nonzero component is the effective cycle of its
maximal-dimensional support with positive generic lengths. A coherent
direct-sum decomposition therefore has only one nonzero summand. For an
equivariant decomposition, S-21's identity-sector Hochschild isomorphism
and invariant CY4 adjunction give exactly the same argument; the common
nonzero scalar normalization by the group order changes no kernel or
image. ∎

This does not exclude an indecomposable nonsimple complex or sheaf.
Together with S-21, it says that a coherent solution in the
imaginary-secant branch must be indecomposable and have endomorphism
dimension greater than two. The nonsimple control
\(\mathcal O_X\oplus\mathcal O_p\) from S-23 cannot itself be turned
into such a character by derived equivalence, since its two summands
have nonzero characters and the argument is categorical.

The same direct-sum conclusion holds for the balanced RM character of
S-22: use (E801.5) in place of (E801.3). Every nonzero summand character
is then proportional to the whole balanced character and has contraction
rank eighteen, whereas the trace image of the sum has dimension
eighteen. Thus a coherent fully semiregular balanced-RM realization is
necessarily indecomposable, with no simplicity or rank hypothesis.

\subsection{SF-E803. Surface point-ideal products cannot become the
genuine RM
character}\label{SF-E803}

Let \(S_1,S_2\) be abelian surfaces, let \(p_i\in S_i\), and put

\[
 X=S_1\times S_2,
 \qquad F=I_{p_1}\boxtimes I_{p_2}.
\]

Although \(F\) is a simple gluable fully semiregular object, no sequence
of derived equivalences of abelian varieties, isogeny pullbacks or
pushforwards, tensor products by line bundles, shifts, and nonzero
scalar multiplication of characters can give a character \(w\) on an
abelian fourfold \(Y\) with

\[
 e^B w=N\left(1-\frac q2\theta_1^2\right)
          \left(1-\frac q2\theta_2^2\right),
 \qquad Nq\ne0,
 \tag{E803.1}
\]

where \(B\) is a rational algebraic divisor and \(\theta_1,\theta_2\)
are the two eigenspace components of a compatible polarization for a
genuine real quadratic field acting with multiplicities ((2,2)). The
obstruction does not depend on the output rank or on an integral
normalizer.

\textbf{Proof of the input properties.} On an abelian surface, \(I_p\)
is simple of rank one and has no negative self-Ext. Serre duality gives
\(\operatorname{ext}^2=1\), and its degree-two trace is nonzero because
the trace of the identity in degree zero is its rank. Riemann--Roch
gives \(\chi(I_p,I_p)=-2\), hence \(\operatorname{ext}^1=4\). Its first
contraction rank is four: the translation directions detect the point
and the degree-zero tensor directions detect the rank, in separate Hodge
components. Thus the degree-one full trace is injective as well. The
product trace is injective on the three Künneth summands in degree two,
lying in distinct pairs of Hodge differences. This proves the assertions
for \(F\), with \(\operatorname{ext}^1(F,F)=8\),
\(\operatorname{ext}^2(F,F)=18\), and \(\chi(F,F)=4\).

\textbf{Proof of the obstruction.} For \(X\), the Fock form of
\(1-[p_i]\) is a nondegenerate two-vector \(\omega_i\) on

\[
 U_i=\operatorname{Lie}(S_i\times\widehat S_i)
       \subset W_X.
\]

Indeed its two nonzero terms are the volume forms of the two-dimensional
translation and tensor directions; their coefficients are nonzero. The
product character is therefore \(\omega_1\wedge\omega_2\). By SF-E801,
its degree-two contraction kernel intrinsically recovers the unordered
pair \(U_1,U_2\). Both recovered spaces are Lie algebras of abelian
subvarieties of \(X\times\widehat X\).

This property is preserved by every listed operation. A derived
equivalence induces an isomorphism of the abelian varieties
\(X\times\widehat X\) and \(Y\times\widehat Y\), whose differential is
the induced map on \(HH^1\); naturality of Hochschild cap product
carries the contraction kernel with it. This is the exact
translation--tensor conjugation assertion of Orlov Theorem2.10 and
Corollary2.13, not an arbitrary complex linear change of coordinates. An
isogeny \(f:Y\to X\) transports \(W_X\) to \(W_Y\) by
\((df^{-1},d\widehat f)\), the differential of a rational isogeny. On
rational cohomology, pushforward differs from inverse pullback by the
nonzero degree scalar. Thus both preserve the property that a recovered
subspace is the Lie algebra of an abelian subvariety. A rational
algebraic B-field is the differential of the rational shear

\[
 (x,\alpha)\longmapsto(x,\alpha+\phi_B(x)),
\]

where \(\phi_B\in\operatorname{Hom}(X,\widehat X)\otimes\mathbf Q\); it
has inverse obtained by replacing \(B\) with (-B). Clearing its
denominator shows that it also preserves algebraic subgroup directions.
Shifts and nonzero character scalars do not change the contraction
kernel.

For (E803.1), on the other hand, SF-E801 recovers the two spaces

\[
 W_\pm=H^0(T_Y)_\pm\oplus H^1(\mathcal O_Y)_\pm,
\]

the eigenspaces of the diagonal real-field action on
\(Y\times\widehat Y\). The polarization identifies the action on the
dual with the same real embedding, since it is Rosati invariant. Neither
\(W_+\) nor \(W_-\) is the Lie algebra of a positive-dimensional abelian
subvariety. Choose a nonrational real field element \(u\). If such a
subvariety existed, its underlying real tangent space and rational first
homology would lie in an eigenspace of the rational operator \(u\) with
irrational real eigenvalue \(\sigma(u)\). A nonzero rational vector
\(z\) cannot satisfy \(uz=\sigma(u)z\): any nonzero rational coordinate
of \(z\) would make \(\sigma(u)\) rational. This is a contradiction.

Thus the intrinsic recovered pair would have to be both algebraic and
nonalgebraic, proving the obstruction. ∎

The same proof applies to an external product of any two surface
characters whose Fock two-vectors are nondegenerate. It excludes a
concrete higher-rank construction strategy beyond the rank-one and
low-rank tests S-30/S-31. It does not exclude a new indecomposable
fourfold object with the required genuine RM character. SF-E805
strengthens the stated point-ideal-product obstruction to allow
interposed flat Hodge deformations; it does not assume that an inverse
equivalence deforms.

\subsection{SF-E804. A necessary Euler inequality for mixed-constituent
extensions}\label{SF-E804}

Let \(\mathcal H\) be a bounded heart in the ordinary or effective
finite equivariant fourfold category. Suppose \(E\in\mathcal H\) has a
finite filtration with constituents \(F_1,\ldots,F_s\), occurring with
multiplicities \(m_i\), such that

\[
 \operatorname{Hom}(F_i,F_j)=0\ (i\ne j),
 \qquad\operatorname{End}(F_i)=\mathbf C.
\]

If \(E\) has a nonzero normalized imaginary-secant character and is
fully semiregular, then

\[
 \chi(E,E)\le2\sum_i m_i^2-4.
 \tag{E804.1}
\]

Here \(\chi\) is the Euler pairing in the category in which the
filtration is taken. In particular, it is impossible that every
\(\chi(F_i,F_i)\ge2\) and every \(\chi(F_i,F_j)\ge0\) for \(i\ne j\).

\textbf{Proof.} Applying left exact Hom successively to both filtrations
gives

\[
 \dim\operatorname{Hom}(E,E)
       \le\sum_{i,j}m_im_j\dim\operatorname{Hom}(F_i,F_j)
       =\sum_i m_i^2.
\]

S-21 gives \(e_2=12\) and \(e_1\ge8\). CY4 duality and gluability give
\(\chi(E,E)=2e_0-2e_1+12\le2e_0-4\), proving (E804.1). Euler bilinearity
gives the final contradiction under the stated diagonal and off-diagonal
conditions. ∎

This extends the numerical heart-filtration test of S-24 to distinct
mutually Hom-orthogonal constituents. A remaining mixed-constituent
proposal must fail at least one of these explicit Euler inequalities or
the Hom-orthogonality hypothesis; this is a necessary gate, not an
existence assertion.

\subsection{SF-E805. The rational stabilizer algebra closes the
deformation
escape}\label{SF-E805}

For a rational even class \(v\) on an abelian fourfold, let

\[
 H=H_1(X,\mathbf Q),\qquad
 V=H\oplus H^*,\qquad S=\bigwedge^*H^*,\qquad
 \mathfrak g_v=\{a\in\mathfrak{so}(V):a\cdot v=0\}.
\]

Here \(x+\xi\in V\) acts on \(S\) by contraction with \(x\) plus wedge
with \(\xi\), and the orthogonal Lie algebra has its resulting spin
action. Define the rational algebra

\[
 \mathcal C(v)=\operatorname{End}_{\mathfrak g_v}(V).
 \tag{E805.1}
\]

For the point-ideal product character \(v=(1-[p_1])(1-[p_2])\),

\[
 \mathcal C(v)\simeq\mathbf Q\times\mathbf Q.
 \tag{E805.2}
\]

For the genuine real-quadratic balanced character in (E803.1), with
field \(F\),

\[
 \mathcal C(e^{-B}Nv_q)\simeq F.
 \tag{E805.3}
\]

These algebras are preserved up to isomorphism by the operations in
SF-E803 and by parallel transport of the full rational character in a
smooth family of abelian varieties. Consequently SF-E803 remains valid
with any number of such flat Hodge deformations interposed among the
listed operations.

\textbf{Explicit stabilizer calculation.} Work first over any
characteristic-zero field \(k\). Let \(H_i\) have dimension four, let
\(\nu_i\in\bigwedge^4H_i^*\) be nonzero, and put

\[
 v_i=1-c_i\nu_i,\qquad c_i\in k^*,\qquad
 V_i=H_i\oplus H_i^*.
\]

The three blocks of \(\mathfrak{so}(V_i)\) are
\(\operatorname{End}(H_i)\), \(\bigwedge^2H_i\), and
\(\bigwedge^2H_i^*\). For \(A\in\operatorname{End}(H_i)\), the spin
action is \(-A^*+\operatorname{tr}(A)/2\), where \(A^*\) denotes the
induced derivation on the exterior algebra. Its degree-zero and
degree-four components on \(v_i\) are

\[
 \frac{\operatorname{tr}(A)}2(1+c_i\nu_i).
\]

The other two blocks act in degree two: contraction by
\(C\in\bigwedge^2H_i\) and wedge by \(b\in\bigwedge^2H_i^*\) give
\(b-c_i\iota_C\nu_i\). Therefore

\[
 \mathfrak g_{v_i}
 =\mathfrak{sl}(H_i)\oplus
  \{C+b:\ b=c_i\iota_C\nu_i\},\qquad
 \dim\mathfrak g_{v_i}=15+6=21.
 \tag{E805.4}
\]

The displayed sum is a vector-space decomposition inside the orthogonal
Lie algebra; its two summands are not being asserted to be ideals.

Its commutant on \(V_i\) is exactly \(k\). The standard and dual
four-dimensional modules for \(\mathfrak{sl}(H_i)\) are nonisomorphic
irreducible modules. Hence an endomorphism commuting with this
subalgebra is scalar \(a\) on \(H_i\) and scalar \(b\) on \(H_i^*\). A
nonzero wedge/contraction pair in (E805.4) has an off-diagonal block, so
commutation forces \(a=b\). These assertions can also be verified
against the elementary matrices of \(\mathfrak{sl}_4\), without a
classification of spin representations.

The Clifford map

\[
 V_i\longrightarrow S_i^{\rm odd},\qquad
 (x,\xi)\longmapsto\xi-c_i\iota_x\nu_i
 \tag{E805.5}
\]

is an isomorphism: its degree-one and degree-three pieces are separate
isomorphisms. For \(V=V_1\oplus V_2\) and \(v=v_1\otimes v_2\), the
mixed block of the orthogonal Lie algebra is \(V_1\otimes V_2\). Its
action on \(v\), up to a fixed nonzero sign, is the tensor product of
(E805.5). It is therefore an isomorphism onto
\(S_1^{\rm odd}\otimes S_2^{\rm odd}\). It cannot cancel the diagonal
blocks, whose images lie in \(S_1^{\rm even}\otimes S_2^{\rm even}\).

For the diagonal blocks, the equation

\[
 a_1v_1\otimes v_2+v_1\otimes a_2v_2=0
\]

forces \(a_iv_i=\lambda_i v_i\), with \(\lambda_1+\lambda_2=0\).
Comparing degree zero and degree four in the preceding formula forces
\(\lambda_i=0\): the coefficients of \(v_i\) have opposite signs, while
those of \(A\cdot v_i\) have the same sign. The degree-two component
then gives (E805.4). Thus

\[
 \mathfrak g_v=\mathfrak g_{v_1}\oplus\mathfrak g_{v_2}.
 \tag{E805.6}
\]

There are no intertwiners between \(V_1\) and \(V_2\): the first Lie
algebra factor acts without a fixed vector or trivial quotient on
\(V_1\) and trivially on \(V_2\), and conversely. The scalar commutant
on each block consequently gives \(\mathcal C(v)=k\oplus k\).

For a product of surfaces these blocks are rational, proving (E805.2).
For a genuine RM character, extend scalars from \(\mathbf Q\) to its
real quadratic field \(F\). The two eigenspaces of \(H\) have dimension
four over \(F\), and \(q\theta_i^2/2\) is a nonzero multiple of a volume
form on each. The calculation gives
\(\mathcal C(v)\otimes F=F\oplus F\). The nontrivial automorphism of
\(F/\mathbf Q\) exchanges the two eigenspaces and their two scalar
idempotents. The fixed algebra is \(\{(a,\bar a):a\in F\}\), isomorphic
to \(F\). This proves (E805.3) before B-normalization; a rational
B-shear conjugates the algebra.

\textbf{Invariance under transformations and deformation.} A nonzero
scalar multiple of \(v\) has the same stabilizer. A rational isometry of
\(V\), together with an intertwiner of its Clifford modules, conjugates
the stabilizer and its commutant. Line-bundle twists and rational
B-fields have the explicit exponential Clifford action. Rational isogeny
pullback has the exterior action of the induced rational linear map on
\(H^*\), and pushforward is a scalar times its inverse.

For derived autoequivalences, the required Clifford compatibility is
Golyshev--Lunts--Orlov, Proposition4.3.7 and Lemma4.3.7.1; the complete
proof was checked in the author version pp17--19. Equivalences between
different abelian varieties follow from the same argument with source
and target distinct, or by this explicit reduction. For
\(\Phi:D^b(X)\to D^b(Y)\), apply the autoequivalence result on
\(X\times Y\) to \((\Phi^{-1}\boxtimes\Phi)\circ\mathrm{swap}\). On
cohomology this swaps the tensor factors through \(\Phi^H\) and its
inverse. Restricting the resulting Clifford identity to the source
factor gives the intertwining identity for \(\Phi^H\). This uses an
actual equivalence at the fibre where it is applied.

In a smooth family, rational homology, its dual, the exterior algebra,
and the Clifford action are local systems. A flat full character
therefore gives a fixed stabilizer and commutant under parallel
transport. No inverse equivalence has to extend across the deformation.
Finally, \(\mathbf Q\times\mathbf Q\) has a nontrivial idempotent, while
\(F\) is a field, so the two algebras cannot be isomorphic. This proves
the obstruction through interposed flat deformations. ∎

For a direct untransformed deformation, there is a shorter check
suggested by the parent audit. The degree-four class \([p_1]+[p_2]\) is
the sum of two transverse rational volume forms on \(H^1(X,\mathbf Q)\).
Its mixed wedge kernel recovers the unordered pair of rational
four-dimensional summands. The connected Hodge circle preserves each
summand. They remain rational Hodge substructures and hence give an
isogeny decomposition into abelian surfaces. The Clifford calculation
handles the transformed setting, where an individual degree-four
component need not stay separately recognizable.

The same deformation obstruction holds for a product of any two
nonzero-rank surface characters with nonzero self-Euler pairings.
Dividing each character by its rank and multiplying by the rational
exponential of minus its slope puts it in the form \(1-c_i\nu_i\). The
self-Euler hypothesis is precisely \(c_i\ne0\). The slope is a rational
algebraic divisor, so this is one of the rational Clifford shears
already treated. Thus using different positive-rank surface constituents
does not evade SF-E805.

\subsection{SF-E806. An all-rank integral obstruction on a principal
real-multiplication
datum}\label{SF-E806}

Let \(A=E_i^4\) with its product principal polarization \(\theta\), and
let

\[
 U=\operatorname{diag}\left(
 \begin{pmatrix}1&1\\1&-1\end{pmatrix},
 \begin{pmatrix}1&1\\1&-1\end{pmatrix}\right).
 \tag{E806.1}
\]

This integral symmetric endomorphism satisfies \(U^2=2I\), so it gives a
Rosati-invariant maximal-order action of \(F=\mathbf Q(\sqrt2)\), of
multiplicity two at each real embedding on the tangent space. Write
\(\theta_\pm\) for the two real eigencomponents of \(\theta\), and put

\[
 C=\frac{\theta_+^2+\theta_-^2}{2}.
\]

\textbf{The integral assertion.} For every positive integer \(r\) and
every integral class \(\ell\in H^2(A,\mathbf Z)\),

\[
 \frac{\ell^2}{2r}-C\notin H^4(A,\mathbf Z).
 \tag{E806.2}
\]

This quantifies over all integral degree-two classes, not only algebraic
classes or the RM divisor plane. The same assertion holds on any
deformation preserving this integral polarized RM datum.

Consequently there is no positive-rank perfect object \(E\), with any
rational divisor normalization \(B\), satisfying

\[
 e^B\operatorname{ch}(E)
 =r\left(1-\frac{\theta_+^2}{2r}\right)
     \left(1-\frac{\theta_-^2}{2r}\right).
 \tag{E806.3}
\]

In particular, the datum has no ordinary simple gluable fully
semiregular object of the balanced RM character in any positive rank.
Indeed S-22 gives \(4q^2r^2\in\{2,4\}\) for such an object. Since \(q\)
is positive rational, the value two would make \((qr)^2=1/2\),
impossible in \(\mathbf Q\). Thus \(q=1/r\), which is (E806.3). This
last implication uses the supplied S-22; the integral assertion below is
independent of semiregularity and simplicity.

\subsubsection{The integral four-form and its obstruction
tensor}\label{the-integral-four-form-and-its-obstruction-tensor}

Choose an integral basis \(a_0,b_0,\ldots,a_3,b_3\) of
\(H^1(A,\mathbf Z)\), with \(\theta=\sum_i a_i b_i\); juxtaposition
means exterior product. Set

\[
 \theta_U=\sum_{i,j}U_{ij}a_i b_j,\qquad
 \theta_\pm=\frac{\theta\pm\theta_U/\sqrt2}{2},\qquad
 Q=4C=\theta^2+\frac{\theta_U^2}{2}.
 \tag{E806.4}
\]

The divided square \(x^{[2]}=x^2/2\) of an integral two-form is
integral. Thus \(Q\) is integral. Make the unimodular integral change of
basis

\[
 \begin{aligned}
 s_0&=a_0+a_1,&s_1&=b_0+b_1,&t_0&=a_1,&t_1&=b_1,\\
 s_2&=a_2+a_3,&s_3&=b_2+b_3,&t_2&=a_3,&t_3&=b_3 .
 \end{aligned}
\]

Let \(S=\langle s_0,\ldots,s_3\rangle\),
\(T=\langle t_0,\ldots,t_3\rangle\), and

\[
 \omega=s_0s_1+s_2s_3,\quad
 \tau=t_0t_1+t_2t_3,\quad
 m=s_0t_1+t_0s_1+s_2t_3+t_2s_3.
\]

Direct substitution gives

\[
 \theta=\omega-m+2\tau,\qquad
 \theta_U=\omega-2\tau.
 \tag{E806.5}
\]

Consequently

\[
 Q\bmod2=\omega^{[2]}=s_0s_1s_2s_3,\qquad
 R:=\frac{Q_{(2,2)}}2\bmod2
    =m^{[2]}+\omega\otimes\tau .
 \tag{E806.6}
\]

The subscript denotes the \(S,T\) bidegree. The mixed two-form \(m\) is
a tensor of rank four in \(S\otimes T\). Its divided square is the
exterior-square tensor of the corresponding isomorphism, and therefore
has rank six in

\[
 \bigwedge^2(S/2S)\otimes\bigwedge^2(T/2T).
\]

Adding \(\omega\otimes\tau\), a rank-one tensor, leaves
\(\operatorname{rank}R\ge5\). This is the only numerical rank needed
below. The exact certificate computes rank six. The first identity in
(E806.6) also shows that \(Q\) is primitive at two.

\subsubsection{A two-adic divided-square lifting
lemma}\label{a-two-adic-divided-square-lifting-lemma}

Let \(L=S\oplus T\) be a free \(\mathbf Z_2\)-module with both summands
of rank four. Suppose \(Q\in\bigwedge^4L\) has reduction equal to a
nonzero volume form of \(S/2S\), and suppose the tensor

\[
 R=\frac{Q_{(2,2)}}2\bmod2
 \in\bigwedge^2(S/2S)\otimes\bigwedge^2(T/2T)
\]

has rank at least two. Then, for no positive integer \(r\) and no
\(\ell\in\bigwedge^2L\), is

\[
 \ell^{[2]}-\frac r4Q\in r\bigwedge^4L.
 \tag{E806.7}
\]

The reduction assumption makes \(Q_{(2,2)}\) divisible by two, so the
definition of \(R\) is meaningful.

\textbf{Proof.} Integrality of the left side forces \(4\mid r\), because
\(Q\) has a unit coefficient. Write \(r=2^e k\), with \(k\) odd and
\(e\ge2\). The alleged congruence is

\[
 \ell^{[2]}\equiv 2^{e-2}kQ\pmod {2^e}.
 \tag{E806.8}
\]

Its left side has minimum coefficient valuation exactly \(e-2\). If
\(\ell=2^a\ell'\) with \(\ell'\) primitive, then \(2a\le e-2\). Dividing
by \(2^{2a}\) reduces to the same congruence for \(\ell'\) and exponent
\(e'=e-2a\ge2\). It therefore suffices to consider primitive \(\ell\).

If \(e=2\), then \(\ell^{[2]}\bmod2\) is the nonzero decomposable volume
form of \(S/2S\). An alternating two-form over a field of characteristic
two whose divided square is a nonzero decomposable four-form has rank
four and support equal to the support of that four-form. Here is a
direct check. In a symplectic basis it is a sum of disjoint pairs. If it
has at least two pairs, every vector in its support occurs nontrivially
in a contraction of its divided square, so the support of the divided
square equals its own support. The nonzero decomposable divided square
has support dimension four. Thus

\[
 \ell=\ell_0+2D,\qquad
 \ell_0\in\bigwedge^2S,
\]

with \(\ell_0\bmod2\) nondegenerate. In bidegree \((2,2)\), modulo four,
its divided square is

\[
 \ell^{[2]}_{(2,2)}
   \equiv 2\ell_0\otimes D_{(0,2)}\pmod4.
\]

After division by two and reduction modulo two, this has tensor rank at
most one. The right side of (E806.8) has tensor rank
\(\operatorname{rank}R\ge2\). This is a contradiction.

Now suppose \(e\ge3\). The alternating elementary-divisor algorithm over
\(\mathbf Z_2\) puts \(\ell\) in the form of a sum of disjoint pairs
with coefficients \(2^{a_1},\ldots,2^{a_j}\), absorbing units into basis
vectors. This algorithm follows by choosing an entry of minimum
valuation, clearing its row and column using divisibility by that entry,
and continuing on the complementary summand. Since \(\ell\) is
primitive, \(a_1=0\). The minimum coefficient valuation of
\(\ell^{[2]}\) is \(a_1+a_2=e-2\), so \(a_2=e-2\) and every subsequent
coefficient has at least that valuation. In the original \(S,T\)
coordinates we may consequently write

\[
 \ell=\alpha+2^{e-2}\beta,
\]

where \(\alpha\) is primitive and decomposable and \(\beta\) is
integral. Dividing (E806.8) by \(2^{e-2}\) gives

\[
 \alpha\beta+2^{e-2}\beta^{[2]}\equiv kQ\pmod4.
 \tag{E806.9}
\]

Let \(P\) be the two-dimensional support of \(\bar\alpha=\alpha\bmod2\).
Reducing (E806.9) modulo two gives

\[
 \bar\alpha\,\bar\beta=\operatorname{vol}_{S/2S}.
 \tag{E806.10}
\]

Every vector in \(P\) wedges to zero with the left side, hence with the
right side. The wedge-annihilator of this nonzero decomposable volume
form is \(S/2S\), so \(P\subset S/2S\). Comparing the \((2,2)\) and
\((3,1)\) components in (E806.10) now gives

\[
 \bar\beta_{(0,2)}=0,\qquad
 \bar\beta_{(1,1)}\in P\otimes(T/2T).
 \tag{E806.11}
\]

Indeed wedge multiplication by the nonzero \(\bar\alpha\in\bigwedge^2P\)
is injective on \(\bigwedge^2(T/2T)\); on \(S/2S\) its kernel is exactly
\(P\).

Write \(\alpha=a b\), with \(\bar a,\bar b\) a basis of \(P\). The
\(T\)-components of \(a,b\) are divisible by two. Thus
\(\alpha_{(0,2)}\) is divisible by four, \(\alpha_{(1,1)}\) is divisible
by two, and

\[
 (\alpha_{(1,1)}/2)\bmod2\in P\otimes(T/2T).
\]

Consider the \((2,2)\) component of \(\alpha\beta\), divide by two, and
reduce modulo two. Its three contributions are

\[
 \alpha_{(2,0)}\beta_{(0,2)},\quad
 \alpha_{(1,1)}\beta_{(1,1)},\quad
 \alpha_{(0,2)}\beta_{(2,0)}.
\]

The first, after division, has its \(S\)-factor in the line
\(\mathbf F_2\bar\alpha\). The second has both of its \(S\)-factors in
\(P\) by (E806.11), so its \(S\)-factor lies in the same line
\(\bigwedge^2P\). The third vanishes after division and reduction.

If \(e\ge4\), the correction \(2^{e-2}\beta^{[2]}\) in (E806.9) vanishes
modulo four. If \(e=3\), dividing its \((2,2)\) component by two gives

\[
 (\bar\beta^{[2]})_{(2,2)}
  =\bar\beta_{(2,0)}\bar\beta_{(0,2)}
     +(\bar\beta_{(1,1)})^{[2]}.
\]

The first term is zero by (E806.11); the second again has its
\(S\)-factor in \(\bigwedge^2P\). In every case the \((2,2)\) component
of the left side of (E806.9), divided by two modulo two, belongs to

\[
 \mathbf F_2\bar\alpha\otimes
       \bigwedge^2(T/2T),
\]

and has tensor rank at most one. The right side gives \(kR=R\), of rank
at least two. This final contradiction proves the lemma. ∎

\subsubsection{Application to Chern
characters}\label{application-to-chern-characters}

Apply the lemma to (E806.6), where \(\operatorname{rank}R=6\). This
proves (E806.2). If an object satisfied (E806.3), its degree-two part
would force \(B=-\ell/r\), where \(\ell=c_1(E)\). Comparing degree four
then gives

\[
 \operatorname{ch}_2(E)=\frac{\ell^2}{2r}-C.
 \tag{E806.12}
\]

The left side is integral on a complex torus: the formula
\(\operatorname{ch}_2(E)=\ell^2/2-c_2(E)\) holds for perfect K-classes,
the second Chern class is integral, and the divided square of every
integral two-form is integral in the exterior cohomology lattice. This
contradicts (E806.2).

The obstruction is topological in the displayed integral RM datum. It
therefore persists when the point is deformed to a simple or non-CM
fourfold with the same datum. It is stronger than the
minimal-divisor-space hypothesis in supplied S-32 for this datum,
because every integral \(\ell\) was allowed. The following local
argument extends the assertion to other integral data without imposing a
hypothesis on the Néron--Severi space.

\subsection{SF-E807. Maximal-order extension at
two}\label{SF-E807}

Let \(F=\mathbf Q(\sqrt{\Delta})\), where \(\Delta>1\) is even and
squarefree. Suppose a polarized abelian fourfold \((X,\theta)\) has
Rosati-invariant RM by \(F\), with tangent multiplicities \((2,2)\), and
suppose its integral homology at two is stable under the maximal local
order

\[
 R=\mathbf Z_2[u],\qquad u^2=\Delta.
\]

Assume that \(D=\int_X\theta^4/24\) is odd. For every
\(\ell\in H^2(X,\mathbf Z_2)\), every positive integer \(r\), and every
\(z\in\mathbf Q_2\) of valuation zero,

\[
 \frac{\ell^2}{2r}-zC\notin H^4(X,\mathbf Z_2),\qquad
 C=\frac{\theta_+^2+\theta_-^2}{2}.
 \tag{E807.1}
\]

Thus no perfect object, with any rational divisor normalization, has
balanced normalized character \(r v_q\) when \(rq\) is a two-adic unit.
No assumption on \(\operatorname{NS}(X)_{\mathbf Q}\) is needed. In
particular every simple gluable fully semiregular lift for an effective
finite coherent target action \(G\) as in S-22 satisfies

\[
 2\mid D|G|.
 \tag{E807.2}
\]

\textbf{Proof of the local lattice normalization.} Let
\(L=H_1(X,\mathbf Z_2)\) and let \(E\) be its alternating polarization
pairing. Since \(D\) is odd, \(E\) is perfect over \(\mathbf Z_2\). The
polynomial \(t^2-\Delta\) is Eisenstein at two. More explicitly, for
\(a,b\in\mathbf Z_2\), the valuation of \(a+bu\) in
\(F\otimes\mathbf Q_2\), normalized by \(v(u)=1\), is

\[
 v(a+bu)=\min\{2v_2(a),\,2v_2(b)+1\};
\]

the two possible valuations have different parity. This gives the
discrete valuation ring \(R\), whose uniformizer is \(u\). The
torsion-free finite \(R\)-module \(L\) is free of rank four, by the
elementary basis algorithm over a discrete valuation ring.

Define an \(R\)-valued form

\[
 \Psi(x,y)=E(x,uy)+uE(x,y).
 \tag{E807.3}
\]

The Rosati-invariance identity \(E(ux,y)=E(x,uy)\) shows directly that
\(\Psi\) is \(R\)-bilinear. It is alternating:
\(E(x,ux)=E(ux,x)=-E(x,ux)\), hence \(E(x,ux)=0\), as the target has
characteristic zero.

Put \(\kappa=1/(2u)\). For \(a+bu\in R\), one has

\[
 \operatorname{Tr}(\kappa(a+bu))=b,\qquad
 \operatorname{Tr}(\kappa u(a+bu))=a.
\]

The trace pairing \((v,w)\mapsto\operatorname{Tr}(\kappa vw)\) is
therefore a perfect \(\mathbf Z_2\)-pairing on \(R\), and

\[
 E(x,y)=\operatorname{Tr}(\kappa\Psi(x,y)).
\]

It follows that \(\Psi:L\longrightarrow\operatorname{Hom}_R(L,R)\) is an
isomorphism: composing with the displayed perfect trace identification
is exactly the given perfect \(\mathbf Z_2\)-duality for \(E\). Choose a
symplectic \(R\)-basis \(x_1,y_1,x_2,y_2\), with
\(\Psi(x_i,y_j)=\delta_{ij}\). Such a basis follows by choosing a
primitive vector \(x_1\), finding \(y_1\) with pairing one using
perfectness, splitting off their span by the usual orthogonal
projection, and repeating.

Let \(\alpha_i,\beta_i,\gamma_i,\delta_i\) be the \(\mathbf Z_2\)-dual
basis to \(x_i,ux_i,y_i,uy_i\). Directly evaluating \(E\) and
\(E(u\,\cdot,\cdot)\) gives

\[
 \theta=m:=\sum_{i=1}^2(\alpha_i\delta_i+\beta_i\gamma_i),
 \qquad
 \theta_u=\omega+\Delta\tau,
 \tag{E807.4}
\]

where

\[
 \omega=\alpha_1\gamma_1+\alpha_2\gamma_2,\qquad
 \tau=\beta_1\delta_1+\beta_2\delta_2.
\]

Take \(S=\langle\alpha_1,\gamma_1,\alpha_2,\gamma_2\rangle\) and
\(T=\langle\beta_1,\delta_1,\beta_2,\delta_2\rangle\). The mixed tensor
\(m\) is again nondegenerate. Write \(\Delta=2h\), with \(h\) odd. The
eigencomponent formula gives

\[
 Q=4C=\theta^2+\frac{\theta_u^2}{\Delta}
   =2m^{[2]}+\frac1h\omega^{[2]}
          +2\omega\tau+4h\tau^{[2]}.
 \tag{E807.5}
\]

This is integral over \(\mathbf Z_2\), its reduction is
\(\omega^{[2]}=\operatorname{vol}_{S/2S}\), and

\[
 (Q_{(2,2)}/2)\bmod2=m^{[2]}+\omega\otimes\tau
\]

has rank six by the same exterior calculation as before. Multiplication
of \(Q\) by the unit \(z\) leaves these two mod-two properties
unchanged. The divided-square lifting lemma applied to \(zQ\) proves
(E807.1).

For the perfect-object assertion, degree-two normalization again forces
\(B=-c_1(E)/r\), and the degree-four comparison gives
\(\operatorname{ch}_2(E)=c_1(E)^2/(2r)-(rq)C\). This contradicts
(E807.1) when \(rq\) is a unit.

Finally suppose \(D|G|\) were odd. S-22 gives

\[
 \frac{4D(rq)^2}{|G|}=k,\qquad k\in\{2,4\}.
\]

For \(k=2\), the valuation of the rational square \((rq)^2\) would be
\(-1\), impossible. For \(k=4\), \(rq\) is a two-adic unit, contrary to
the integral assertion. This proves (E807.2). ∎

The extension allows every integral degree-two class, including extra
classes at CM points. It covers every maximal-at-two lattice for even
squarefree \(\Delta\), and every odd-degree compatible polarization.
SF-E808 below treats the odd squarefree radicands ramified at two. A
nonmaximal order at two and an even value of \(D|G|\) remain outside the
combined obstruction. No eightfold supplier is constructed.

\subsection{SF-E808. Odd radicands ramified at two: an unrestricted
Chern-class
obstruction}\label{SF-E808}

Let \(F=\mathbf Q(\sqrt{\Delta})\), with \(\Delta>1\) squarefree and
\(\Delta\equiv3\pmod4\). Suppose a polarized abelian fourfold
\((X,\theta)\) has Rosati-invariant RM by \(F\), with tangent
multiplicities \((2,2)\), and that \(L=H_1(X,\mathbf Z_2)\) is stable
under the maximal local order \(R=\mathbf Z_2[u]\), \(u^2=\Delta\).
Suppose \(D=\int_X\theta^4/24\) is odd. Define

\[
 C=\frac{\theta_+^2+\theta_-^2}{2}.
\]

\textbf{Assertion.} For every integral class
\(\ell\in H^2(X,\mathbf Z_2)\), every positive integer \(r\), and every
two-adic unit \(z\in\mathbf Q_2\),

\[
 \frac{\ell^2}{2r}-zC\notin H^4(X,\mathbf Z_2).
 \tag{E808.1}
\]

Consequently no positive-rank perfect object, with any rational divisor
normalization, has normalized balanced character \(r v_q\) when \(rq\)
is a two-adic unit. A simple gluable fully semiregular lift with an
effective finite coherent target action \(G\) as in supplied S-22
therefore satisfies \(2\mid D|G|\).

Together with SF-E807, this removes the minimal-divisor-space hypothesis
from supplied S-32 for all real quadratic fields ramified at two.
Maximal-order stability at two and oddness of the relevant degree
product are retained.

\subsubsection{The local polarized RM
lattice}\label{the-local-polarized-rm-lattice}

The element \(u-1\) satisfies

\[
 t^2+2t+1-\Delta=0.
\]

This polynomial is Eisenstein at two, because \(1-\Delta\) has valuation
one. Thus \(R\) is a discrete valuation ring with uniformizer
\(\pi=u-1\). Concretely \(R=\mathbf Z_2+\mathbf Z_2\pi\), and the
valuation of \(a+b\pi\) is \(\min\{2v_2(a),2v_2(b)+1\}\); the two terms
have distinct valuation parity. This identifies \(R\) with the valuation
ring. The finite torsion-free \(R\)-module \(L\) is free of rank four.

Let \(E\) be the polarization pairing on \(L\). Its determinant is odd,
so it is perfect over \(\mathbf Z_2\). Exactly as in the full local
proof of SF-E807, the form

\[
 \Psi(x,y)=E(x,uy)+uE(x,y)
\]

is alternating and \(R\)-bilinear, and it is perfect because

\[
 E(x,y)=\operatorname{Tr}_{(F\otimes\mathbf Q_2)/\mathbf Q_2}
              \left(\frac{\Psi(x,y)}{2u}\right).
\]

For completeness, the trace pairing on \(R\) is perfect: it sends
\(a+bu\) to \(b\), and multiplication first by \(u\) sends it to \(a\).
This identifies the \(R\)-linear dual with the \(\mathbf Z_2\)-linear
dual. A symplectic \(R\)-basis \(x_i,y_i\), \(i=1,2\), follows by
splitting off pairs with pairing one.

Let \(\alpha_i,\beta_i,\gamma_i,\delta_i\) be the dual integral basis to
\(x_i,ux_i,y_i,uy_i\). Put

\[
 \begin{aligned}
 S&=\langle\alpha_1,\gamma_1,\alpha_2,\gamma_2\rangle,
 &T&=\langle\beta_1,\delta_1,\beta_2,\delta_2\rangle,\\
 \omega&=\alpha_1\gamma_1+\alpha_2\gamma_2,
 &\tau&=\beta_1\delta_1+\beta_2\delta_2,\\
 m&=\alpha_1\delta_1+\beta_1\gamma_1
       +\alpha_2\delta_2+\beta_2\gamma_2.
 \end{aligned}
\]

Evaluation of the pairings gives \(\theta=m\) and
\(\theta_u=\omega+\Delta\tau\). Hence

\[
 Q:=2C=m^{[2]}+\frac1\Delta\omega^{[2]}
            +\omega\tau+\Delta\tau^{[2]}.
 \tag{E808.2}
\]

This is an integral two-adic four-form. Its reduction is independent of
the odd radicand:

\[
 \bar Q=m^{[2]}+\omega^{[2]}+\omega\tau+\tau^{[2]}.
 \tag{E808.3}
\]

Here the displayed operations are divided powers and exterior products
reduced from the integral lattice; signs disappear modulo two.

\subsubsection{A contraction with no wedge
annihilator}\label{a-contraction-with-no-wedge-annihilator}

Write \(s_0=\alpha_1,s_1=\gamma_1,s_2=\alpha_2,s_3=\gamma_2\), and
\(t_0=\beta_1,t_1=\delta_1,t_2=\beta_2,t_3=\delta_2\). Let \(\lambda\)
be the dual functional selecting \(s_0\). Direct contraction of (E808.3)
gives

\[
 \chi=\iota_\lambda\bar Q
 =s_1s_2s_3+s_1t_2t_3+s_2t_1t_3+s_3t_1t_2.
 \tag{E808.4}
\]

The map

\[
 \bar L^*\longrightarrow\bigwedge^4\bar L^*,\qquad
 v\longmapsto v\chi
 \tag{E808.5}
\]

is injective. This is a direct monomial calculation, not a generic-rank
assertion. Rename the six variables appearing in \(\chi\) as
\(a,b,c,d,e,f\), so \(\chi=abc+aef+bdf+cde\). Their six exterior
products are

\[
 \begin{array}{c|l}
 a&abdf+acde\\
 b&abef+bcde\\
 c&acef+bcdf\\
 d&abcd+adef\\
 e&abce+bdef\\
 f&abcf+cdef .
 \end{array}
\]

The twelve monomials are distinct, proving independence of these six
images. Exterior multiplication by either of the remaining variables
\(s_0,t_0\) gives a nonzero form containing that variable. These two
images and the preceding six belong to separate summands according to
their \(s_0,t_0\) content. This proves injectivity on the entire
eight-dimensional space.

In particular \(\bar Q\ne0\), so \(Q\) is primitive at two.

\subsubsection{A divided-square lemma with denominator
two}\label{a-divided-square-lemma-with-denominator-two}

Let \(M\) be a free \(\mathbf Z_2\)-module, let \(Q\in\bigwedge^4M\),
and suppose some functional \(\lambda\in(M/2M)^*\) has the property that
exterior multiplication by \(\chi=\iota_\lambda\bar Q\) is injective on
\(M/2M\). Then there are no \(\ell\in\bigwedge^2M\) and positive integer
\(r\) with

\[
 \ell^{[2]}-\frac r2Q\in r\bigwedge^4M.
 \tag{E808.6}
\]

\textbf{Proof.} The hypothesis makes \(Q\) primitive, so integrality
forces \(r\) even. Write \(r=2^e k\), with \(k\) odd and \(e\ge1\). The
alleged congruence is

\[
 \ell^{[2]}\equiv 2^{e-1}kQ\pmod{2^e}.
 \tag{E808.7}
\]

The minimum valuation of its left side is exactly \(e-1\). Removing the
largest common factor \(2^a\) from \(\ell\) reduces the congruence to
the same one for a primitive two-form and exponent \(e'=e-2a\ge1\). We
may assume that \(\ell\) is primitive.

For \(e=1\), the congruence says \(\bar\ell^{[2]}=\bar Q\). The
contraction identity for divided squares gives

\[
 \chi=(\iota_\lambda\bar\ell)\bar\ell.
\]

If \(\iota_\lambda\bar\ell=0\), this contradicts \(\chi\ne0\). If it is
nonzero, it wedges to zero with \(\chi\), contradicting the assumed
injectivity. Thus \(e=1\) is impossible.

For \(e\ge2\), the alternating elementary-divisor algorithm gives

\[
 \ell=\alpha+2^{e-1}\beta,
\]

where \(\alpha\) is primitive and decomposable. Indeed primitivity makes
the first elementary divisor a unit, while the valuation of the divided
square makes the second elementary divisor have valuation \(e-1\); all
subsequent divisors have at least that valuation. Dividing (E808.7) by
\(2^{e-1}\) and reducing modulo two gives

\[
 \bar Q=\bar\alpha\bar\beta.
\]

The two-dimensional support \(P\) of \(\bar\alpha\) therefore wedges to
zero with \(\bar Q\). There is a nonzero \(v\in P\cap\ker\lambda\).
Contracting \(v\bar Q=0\) by \(\lambda\) gives

\[
 0=\lambda(v)\bar Q+v\iota_\lambda\bar Q=v\chi,
\]

which again contradicts injectivity. Both cases are impossible. ∎

\subsubsection{Chern-character consequence and
scope}\label{chern-character-consequence-and-scope}

Multiplication by the unit \(z\) in (E808.1) leaves \(\bar Q\) and its
contraction unchanged. The lemma applied to \(zQ\), with \(Q=2C\),
proves (E808.1). For a normalized balanced character, degree two forces
\(B=-c_1(E)/r\), while degree four gives

\[
 \operatorname{ch}_2(E)=\frac{c_1(E)^2}{2r}-(rq)C.
\]

The Chern character on the left is integral on a complex torus, since
\(\operatorname{ch}_2=c_1^2/2-c_2\) and every integral two-form has an
integral divided square. Hence \(rq\) cannot be a two-adic unit.

If a simple gluable fully semiregular equivariant lift had \(D|G|\) odd,
supplied S-22 would give \(4D(rq)^2/|G|\in\{2,4\}\). The value two has
impossible odd valuation for a rational square; the value four makes
\(rq\) a two-adic unit. Thus \(2\mid D|G|\).

No restriction to the RM divisor plane was made: every integral
\(H^2\)-class is included, so extra CM divisor classes cannot evade the
conclusion. This does not exclude even-degree polarizations, even-order
effective finite actions, nonmaximal local orders, or other proposed
characters. It constructs no eightfold supplier.

\subsubsection{Source and check status}\label{source-and-check-status}

The proof uses only the supplied S-22 Euler restriction for the final
semiregularity corollary. The integral obstruction, local pairing
calculation, contraction matrix, and lifting argument are proved in full
above. S-32 is a comparison result, not an input to this proof. The
finite certificate checks exact normal-form substitutions and the
contraction rank; it does not replace the proof for every rank or every
radicand.

\subsection{Sources, checks, and remaining
scope}\label{sources-checks-and-remaining-scope}

\begin{itemize}
\item
  DEG\_L S-21--S-25 and S-30--S-34 are supplied background. Their CY4
  adjunction, identity-sector assertion and previously proved
  obstructions are retained with their existing status.
\item
  Orlov, \emph{Derived categories of coherent sheaves on Abelian
  varieties and equivalences between them}, Izv. Math.66(3),569--594
  (2002), Theorem2.10 and Corollary2.13. Publication verified at
  https://www.mathnet.ru/eng/im389 . Full theorem proof and corollary
  proof read in the author manuscript
  https://homepage.mi-ras.ru/\textasciitilde orlov/papers/abelnew.pdf ,
  pp9--12. The publisher PDF was located but its subsequent fetch
  failed; no unverified printed-page locator is asserted. This published
  result supplies the abelian-group transport, while the kernel
  reconstruction above is a new proof.
\item
  \texttt{eightfold\_kernel\_certificate.py} is an exact rational
  exterior-algebra check with its scope stated in its docstring and
  JSON. It is not a proof assistant or a sheaf-construction certificate.
\item
  Golyshev--Lunts--Orlov, \emph{Mirror symmetry for abelian varieties},
  J. Algebraic Geom.10(3),433--496 (2001), Proposition4.3.7 and
  Lemma4.3.7.1. Published work; complete relevant proof inspected in
  https://homepage.mi-ras.ru/\textasciitilde orlov/papers/GLO.pdf ,
  author-version pp17--19. The exact cohomological Clifford intertwining
  is the input used in SF-E805. Printed theorem pagination was not
  independently verified.
\item
  The new integral assertions SF-E806--SF-E808 are proved by elementary
  exterior algebra and the local alternating-module algorithms
  explicitly described above. Their semiregularity corollaries use
  supplied DEG\_L S-22, equations (S22.2)--(S22.3). The comparison with
  S-32 respects that theorem's stated minimal-divisor-space hypothesis.
\item
  The exact arithmetic certificate for SF-E806 checks the integral
  change of coordinates, the primitive four-form, the rank-six
  obstruction tensor, and all twenty-eight mod-two nondegenerate
  alternating forms in the first lifting case. The proof for arbitrarily
  large ranks is the two-adic argument above, not a bounded computation.
\item
  No split-eightfold CM input is constructed here. The original nonsplit
  sixfold gate remains open, narrowed by SF-E802--E808. In particular,
  the success of SF-QC01--05 for split sixfolds does not imply this
  remaining eightfold supplier.
\end{itemize}
